\section{Testing and Experimental Evidence}
\subsection{Environment and test outcomes}

The team recorded the final audit on \textbf{November 13, 2026} with Windows, Python 3.14.5, \texttt{cryptography} 46.0.7, Node.js for JS checks, and Npcap/TShark for loopback capture. \Cref{tab:testresults} distinguishes automated results from browser/capture observations; times are single recorded runs, not performance guarantees.

\begin{table}[H]
\centering\small
\begin{tabularx}{\linewidth}{@{}L{3.3cm}L{2.5cm}Y@{}}
\toprule
\textbf{Test} & \textbf{Recorded result} & \textbf{Verified scope} \\
\midrule
Python unit/integration & \textbf{65/65 PASS} (43.045 s) & TCP/TLS, UDP loss/reorder, registration, outbox, file integrity and cleanup. \\
Smoke test & \textbf{PASS} (1.044 s) & Two-client chat, online list, UDP binary file, hash and receipt. \\
Load test & \textbf{PASS} (20 clients) & 20 simultaneous localhost connections and private chat. \\
Extended concurrency & \textbf{PASS} (2/10/20) & Broadcast/private chat, 1/5/10 concurrent files and resource cleanup. \\
JavaScript logic & \textbf{PASS} (0.213 s) & Registration/login, UI states, progress and message de-duplication (DOM double). \\
Real HTTPS Chrome & \textbf{PASS} (recorded) & Register, login/logout, broadcast, empty/text/binary files, SHA-256 and reconnect/replay. \\
Two separate LAN hosts & \textbf{NOT TESTED} & Tests above were localhost only. \\
\bottomrule
\end{tabularx}
\caption{Final verification results from the committed audit logs.}
\label{tab:testresults}
\end{table}

\subsection{Screenshots of the integrated application}

The screenshots below were captured from the running application and genuine terminal output on November 13, 2026 at commit \texttt{cfe60cd}. Historical Member 4 UI images were not substituted. The separate screenshot run repeated all four requested commands: 65 Python tests passed in 44.532 seconds (45.398 seconds including command overhead); smoke, 20-client load and JavaScript checks all exited successfully (1.216, 2.168 and 0.275 seconds respectively). The earlier audit timings in \Cref{tab:testresults} refer to their original logs.

\begin{figure}[H]
\centering
\evidenceslot{01_https_login_roster.png}{HTTPS login and online roster}{Real trusted Chrome UI with both synthetic accounts connected.}{12cm}
\caption{Actual Chrome account-mode interface: Alice is connected, both users appear in the roster, and chat/file controls are available. The browser address and controls are cropped from the same original screenshot; blank feed space is omitted.}
\label{fig:browser-evidence}
\end{figure}

\begin{figure}[H]
\centering
\evidenceslot{02_chat_file_verification.png}{Messaging and verified file}{Actual sender receipt and recipient SHA-256 verification.}{13cm}
\caption{Actual Alice and Bob browser crops: private and broadcast messages each appear once at Bob, Alice shows DELIVERED, and the 82-byte synthetic file has a successful FILE\_RECEIPT. Source and received SHA-256 match independently.}
\label{fig:file-browser-evidence}
\end{figure}

\begin{figure}[H]
\centering
\evidenceslot{06_tests_and_concurrency.png}{Terminal regression and concurrency evidence}{Genuine terminal output from this screenshot run.}{6cm}
\caption{Crops of actual terminal output from two phases of the same test runner: 65 tests OK and 20 concurrent clients PASS. Complete logs and full-window originals are retained separately.}
\label{fig:terminal-evidence}
\end{figure}

\subsection{Wireshark evidence and integrity checks}

The audit capture \texttt{final-loopback.pcapng}, saved in the dated evidence directory, contains \textbf{389 frames}: 178 TCP and 211 UDP; 70 were decoded as TLS in TShark analysis. A deliberately lost transport ACK resulted in byte-identical UDP DATA retransmissions in \textbf{frames 109 and 120}, followed by an ACK in \textbf{frame 121}; the packet UUID and sequence number did not change. The capture also contains the 43-packet 32-KiB file stream. The capture used \texttt{19100/TCP} and \texttt{19101/UDP}; application defaults are \texttt{9100/9101}.

\begin{table}[H]
\centering\small
\begin{tabularx}{\linewidth}{@{}L{3.6cm}Y@{}}
\toprule
\textbf{Captured observation} & \textbf{Evidence / meaning} \\
\midrule
TCP connection and TLS & SYN/SYN-ACK plus encrypted control sessions; plaintext passwords are not visible. \\
Custom UDP DATA/ACK & Header \texttt{COLB}; DATA kind 1, ACK kind 2, distinct packet IDs and stream IDs. \\
Lost ACK and replay & Frames 109/120 repeat the same application DATA; ACK follows at frame 121. \\
File transfer & Consecutive file-stream sequences 0--42; successful SHA-256 is confirmed by client tests, not inferred from ciphertext. \\
\bottomrule
\end{tabularx}
\caption{Concrete observations from the preserved packet-capture analysis.}
\label{tab:pcap}
\end{table}

\begin{figure}[H]
\centering
\evidenceslot{05_wireshark_udp_retransmission.png}{Wireshark DATA, retransmission and ACK}{Capture on TCP 19100 / UDP 19101; highlight equal DATA frames 109 and 120 and ACK frame 121.}{14cm}
\caption{Actual Wireshark crops of the preserved 389-frame audit trace. DATA frames 109 and 120 have identical payloads; frame 121 is the following ACK. In each byte pane, the custom header starts at packet offset 0x20: kind at 0x25 (01 DATA, 02 ACK), sequence at 0x26--0x29 (zero), and packet UUID at 0x2a--0x39. The matching UUID is 82bc661c-2bf6-4a47-83e5-490095401c5a. TShark independently verified these bytes; all originals are retained.}
\label{fig:wireshark-evidence}
\end{figure}

The trusted-browser account checks independently verified matching SHA-256 for an empty file, a Unicode text file, and a 2-MiB binary file. In the 2-MiB case, a chat message appeared while reception was at \textbf{35\%}, before transfer completion. These results demonstrate stream separation and file integrity under the recorded localhost conditions. The PCAP, Chrome records and terminal logs are preserved in \filename{docs/evidence/} in the repository; they are separate test runs and are not presented as one combined capture.
